%%%%%%%%%%%%%%%%%%%%%%%%%%%%%%%%%%%%%%%%%%%%%%%%%%%%%%%%
\documentclass[10pt, aspectratio=169]{beamer}
\usepackage[many]{tcolorbox}
\usetheme{metropolis}
\usepackage[english]{babel}
\usepackage[utf8]{inputenc}
\usepackage[T1]{fontenc}
\usepackage{array}
\usepackage{booktabs}
\usepackage{colortbl}
\usepackage{multirow}
\usepackage{graphicx}
\usepackage{mathrsfs}
\usepackage{bbding}
\usepackage{pifont}
\usepackage{wasysym}
\usepackage{amssymb}
\usepackage{enumerate}
\usepackage{booktabs,threeparttable}
\usepackage{adjustbox}
\usepackage{tikz}
\usetikzlibrary{calc}
\usepackage{geometry}
\usepackage{appendixnumberbeamer}


% DEFINE COLORPALETTE
\newcommand{\cmark}{\ding{51}}%
\newcommand{\xmark}{\ding{55}}
\definecolor{unibasmintlight}{RGB}{233, 255, 253}
\definecolor{unibasmint}{RGB}{165,215,210}
\definecolor{unibasmintmid}{RGB}{30,165,165}
\definecolor{unibasanthr}{RGB}{45,55,60}
\definecolor{unibasgreydark}{RGB}{140,145,150}
\definecolor{unibasred}{RGB}{210,5,55}
\definecolor{unibaspink}{RGB}{235,130,155}
\definecolor{unibasgrey}{RGB}{190,195,200}
\definecolor{unibasmintdark}{RGB}{0,110,110}
\definecolor{unibasblack}{RGB}{0,0,0}
\definecolor{unibaswhite}{RGB}{255,255,255}

\definecolor{scenA_table}{RGB}{245, 245, 180}
\definecolor{scenB_table}{RGB}{235, 215, 245}
\definecolor{scenC_table}{RGB}{200, 225, 240}

\definecolor{capacity_upgrade}{RGB}{64, 153, 153}
\definecolor{capacity_expansion}{RGB}{162, 216, 216}


% -------------------------------------------------------
% METROPOLIS COLOR SETUP
% -------------------------------------------------------
% \setbeamercolor{normal text}{fg=unibasanthr, bg=white}
% \setbeamercolor{alerted text}{fg=unibasred}
% \setbeamercolor{example text}{fg=unibasmintdark}

% \setbeamercolor{frametitle}{bg=unibasmintlight, fg=unibasmintdark}
% \setbeamercolor{title separator}{fg=unibasmintdark}
% \setbeamercolor{progress bar}{fg=unibasmintdark, bg=unibasgrey}

% \setbeamercolor{block title}{bg=unibasmintdark, fg=unibasmint}
% \setbeamercolor{block body}{bg=unibasmintlight, fg=unibasanthr}
% \setbeamercolor{block title alerted}{bg=unibasred, fg=white}
% \setbeamercolor{block body alerted}{bg=unibasmintlight!50, fg=unibasanthr}

% \setbeamercolor{title}{fg=unibasmintdark}
% \setbeamercolor{subtitle}{fg=unibasgreydark}
% \setbeamercolor{author}{fg=unibasanthr}
% \setbeamercolor{institute}{fg=unibasgreydark}
% \setbeamercolor{date}{fg=unibasgreydark}

% \setbeamercolor{section in toc}{fg=unibasmintdark}
% \setbeamercolor{subsection in toc}{fg=unibasgreydark}

% \setbeamercolor{itemize item}{fg=unibasmintdark}
% \setbeamercolor{itemize subitem}{fg=unibasmintmid}
% \setbeamercolor{itemize subsubitem}{fg=unibasgreydark}
% \setbeamercolor{enumerate item}{fg=unibasmintdark}

% \setbeamercolor{button}{bg=unibasmint, fg=unibasanthr}

% footer / standout
\setbeamercolor{footline}{fg=unibasgreydark, bg=white}

% -------------------------------------------------------
% METROPOLIS FONT / STYLE OPTIONS
% -------------------------------------------------------
\metroset{
  titleformat        = regular,
  titleformat frame  = regular,
  sectionpage        = none,
  numbering          = fraction,
  progressbar        = frametitle,
  block              = fill
}

\usefonttheme[onlymath]{serif}
\setbeamertemplate{caption}[numbered]
\setbeamertemplate{frame footer}{\textcolor{unibasgreydark}{Hydro and Wind Power: Complements or Substitutes? (Discussion)}}


\usepackage[backend=biber, style=apa, maxbibnames=999, minbibnames=1,
            maxcitenames=1, mincitenames=1, uniquelist=true]{biblatex}
\addbibresource{optimalpv_conf.bib}
\AtBeginBibliography{\small}
% \hypersetup{colorlinks=true,
%             citecolor=unibasmintmid,
%             linkcolor=unibasmintdark,
%             urlcolor=unibasmintdark}

% -------------------------------------------------------
% BACKUP FRAME COUNTER
% -------------------------------------------------------
\newcommand{\backupbegin}{
  \newcounter{framenumberappendix}
  \setcounter{framenumberappendix}{\value{framenumber}}
}
\newcommand{\backupend}{
  \addtocounter{framenumberappendix}{-\value{framenumber}}
  \addtocounter{framenumber}{\value{framenumberappendix}}
}

% -------------------------------------------------------
% TITLE INFO
% -------------------------------------------------------
\title[]{Hydro and Wind Power: Complements or Substitutes?∗}
\subtitle{Paper Discussion}
\author{Authors: Gautam Gowrisankaran, Stanley Reynolds, \textbf{Rauli Svento}
\\Discussant: Raul Hochuli}
\institute{WCERE - Portugal}
\date{July 2026}
% \titlegraphic{\includegraphics[width=2.3cm]{UniBas_Logo.jpg}}


%*******************************************************************************
\begin{document}
%*******************************************************************************

% -------------------------------------------------------
\begin{frame}
  \maketitle
\end{frame}


%%%%%%%%%%%%%%%%%%%%%%%%%%%%%%%%%%%%%%%%%%%%%%%
\section{Discussion}

% -------------------------------------------------------
\begin{frame}{Summary}
\begin{itemize}
\item \textbf{Dynamic stochastic model} of a Finnish system operator, \textbf{schedules and dispatches generation} hour-by-hour to maximize expected social surplus.
\item \textbf{Hourly two-stage decision process}
\begin{itemize}
    \item First: Forecasting of load, wind, CHP and hydro inflow $\rightarrow$ schedule fossil fuel generators.
    \item Second: After realization of these key variables $\rightarrow$ dispatch hydro and fossil capacity.
    \item Setup with three generation sources: Thermal plants (merit-order), \textbf{wind} (exogenous capacity and weather-driven capacity factor), \textbf{hydro} (flexible reservoir stock and RoR). 
\end{itemize}
\item \textbf{Weather variables} and \textbf{load/generation} estimated in SUR system, hourly 2022 Finish data.
% \vspace{4pt}
% \item Solves the operator's dynamic programming problem computationally (Halton draws, Markov weather transitions, value function iteration); simulates 100-year paths to capture long-run reservoir variability
% \vspace{4pt}
\item Dynamic model, adjustments in two "dimension"
\begin{itemize}
    \item Scaling up wind capacity by +50\% and +100\% (of 5'000 MW)
    \item +/-10\% and +/-30\% rescale of hydro 
    \item Asses social benefits (operating cost, fossil reserves, price, outage probability, CO2 emissions)
\end{itemize}
\end{itemize}
\end{frame}


% -------------------------------------------------------
\begin{frame}{Main Comments}
\begin{itemize}
    \item Rescaling hydro \textbf{including run-of-river} by adjusting must-run nuclear $\rightarrow$ substitute a flexible for an inflexible resource. 
    \begin{itemize}
        \item \textit{Higher RoR generation exacerbates wind curtailment in spring/summer. }
    \end{itemize}
    \item \textbf{Exclusion of Nord Pool} prohibits export of excess wind to neighbors. 
    \item Using \textbf{2022 as a single weather year}
    \begin{itemize}
        \item \textit{Other years as robustness check.}
        \item \textit{Estimating 8 equation on 8760 hours, might not include 100 years weather variability.} 
    \end{itemize}
    \vspace{0.0cm}
    \item Ext: Assumed \textbf{no price / curtailment response} of consumers 
    % \item Ext: Assumed \textbf{no price / curtailment response} of consumers (contrast to Gowrisankaran et al. 2016; \textit{EV 48\% of new registrations in 2026 (3rd in }Europe)\footnote{https://theicct.org/publication/european-car-market-monitor-may-2026/})
    \begin{itemize}
        \item Contrast to Gowrisankaran et al. 2016; EV 48\% of new Finish cars in \href{https://theicct.org/publication/european-car-market-monitor-may-2026/}{2026}.
        % \item In contrast to Gowrisankaran et al. 2016. 
        % \item \textit{EV 48\% of new registrations in 2026 (3rd in }Europe)\footnote{https://theicct.org/publication/european-car-market-monitor-may-2026/}
        \item \textit{One aggregated battery / EV component (similar to aggregated reservoir)}
        \item \textit{Might allow for higher wind penetration.}
    \end{itemize}
    
    \item Ext: \textbf{Data based (but conservative) emission price} of 80 EUR/tCO2. 
    \begin{itemize}
        \item \textit{Sensitivity analysis for larger range in social cost of carbon . }
        \item \textit{Higher wind share more beneficial, push substitute threshold further upwards.}
    \end{itemize}
    % \item you have investment cost for wind but keep it veery simple with 50 and 100 percent wind icrease - why not optimize for optimal wind capacity?
    % \item model a 100 year model -> no assumptions about temperature adjustments over that time. 
    % \item add model runs with higher soc carbon to see how it might look like with future co2 prices. 
\end{itemize}
\end{frame}

% -------------------------------------------------------
\begin{frame}{Minor Comments}
\begin{itemize}
    \item More runs to show smaller change increments (for hydro and wind)
    \item Converting outcome values (e.g. Table 5, p34) to a line graph (serves complement / substitute story + show "optimal" state).
    \item Single hydro reservoir but explicit mention of individual reservoir specific multiplier (p18).
    \item The paper refers to wind capacity increases of 50\% and 100\% without clearly stating the baseline capacity (5'000\,MW) at first mention, making the interpretation of the counterfactuals harder at a first glance. 
\end{itemize}
\end{frame}

% -------------------------------------------------------
% \begin{frame}{Petty Comments}
% \begin{itemize}
% \end{itemize}
% \end{frame}

\end{document}
